% Generated by prepare_content.py from the accepted Markdown report.
\chapter{MỞ ĐẦU}

\section{Lý do xây dựng website học liệu}

Việc học tập môn Lịch sử Đảng Cộng sản Việt Nam không chỉ đòi hỏi ghi nhớ các mốc thời gian mà còn cần hiểu mối quan hệ giữa bối cảnh lịch sử, đường lối lãnh đạo và kết quả thực tiễn. Đối với giai đoạn kháng chiến chống Mỹ, cứu nước từ năm 1954 đến năm 1975, lượng kiến thức trải dài qua nhiều thời kỳ, gắn với những chuyển biến về nhiệm vụ chiến lược và phương pháp đấu tranh. Điều này đặt ra nhu cầu tổ chức kiến thức thành một hệ thống rõ ràng, thuận tiện cho việc tìm hiểu, đối chiếu và ôn tập.

Từ nhu cầu đó, nhóm lựa chọn xây dựng website học liệu với chủ đề “Từ đường lối đến thực tiễn”. Website là sản phẩm chính của đề tài, kết hợp nội dung lịch sử theo năm giai đoạn với dòng thời gian tương tác, tư liệu minh họa và câu hỏi trắc nghiệm. Thay vì tiếp cận các sự kiện như những thông tin riêng lẻ, người học có thể theo dõi tiến trình, tìm hiểu từng giai đoạn và tự kiểm tra kiến thức trên cùng một giao diện.

Việc liên hệ với chứng tích tại Bảo tàng Chứng tích Chiến tranh bổ sung góc nhìn trực quan cho học liệu. Những hiện vật, hình ảnh và nội dung trưng bày giúp kết nối kiến thức trong giáo trình với điều kiện chiến đấu, đời sống nhân dân và hậu quả của chiến tranh. Hướng tiếp cận này góp phần làm rõ ý nghĩa của các chủ trương lịch sử, đồng thời bồi dưỡng thái độ trân trọng hòa bình và trách nhiệm của thế hệ trẻ.

Là sinh viên Trường Đại học Khoa học Tự nhiên, nhóm mong muốn vận dụng khả năng ứng dụng công nghệ thông tin để tạo ra một công cụ hỗ trợ học tập phù hợp với học phần. Sản phẩm hướng tới việc trình bày kiến thức có chọn lọc, dễ theo dõi và có sự tương tác của người sử dụng. Giáo trình giữ vai trò nền tảng học thuật; website hỗ trợ hệ thống hóa, minh họa và ôn tập, không thay thế việc nghiên cứu tài liệu chính của môn học.

\section{Mục tiêu và nhiệm vụ xây dựng sản phẩm}

\textbf{Mục tiêu tổng quát} là xây dựng một website học liệu tương tác về sự lãnh đạo của Đảng trong cuộc kháng chiến chống Mỹ, cứu nước giai đoạn 1954 - 1975, phục vụ việc tìm hiểu và ôn tập của sinh viên. Nội dung được tổ chức theo hướng kết nối đường lối với thực tiễn, kết hợp kiến thức từ giáo trình với tư liệu chứng tích lịch sử.

\textbf{Mục tiêu cụ thể} gồm ba phương diện. Về nội dung, website hệ thống hóa những kiến thức trọng tâm, làm rõ sự phát triển của đường lối qua từng giai đoạn. Về cách tiếp cận, sản phẩm tạo điều kiện để người học chủ động lựa chọn nội dung, quan sát tư liệu và kiểm tra mức độ ghi nhớ. Về giáo dục, học liệu hướng tới việc bồi dưỡng nhận thức lịch sử, lòng biết ơn và ý thức trách nhiệm đối với cộng đồng, đất nước.

Để thực hiện các mục tiêu trên, nhóm xác định những nhiệm vụ sau:

\begin{enumerate}
\item Hệ thống hóa nội dung theo năm giai đoạn: 1954 - 1960, 1961 - 1965, 1965 - 1968, 1969 - 1973 và 1973 - 1975; mỗi giai đoạn làm rõ bối cảnh, chủ trương của Đảng, kết quả thực tiễn và kiến thức cần ghi nhớ.
\item Tổ chức nội dung dưới dạng trình chiếu kết hợp dòng thời gian tương tác, cho phép người học chuyển giữa các giai đoạn và theo dõi tiến trình lịch sử.
\item Xây dựng khu vực tư liệu chứng tích, gắn hình ảnh và phần giới thiệu với giai đoạn liên quan; gắn tư liệu với thời kỳ tương ứng và hỗ trợ xem thông tin chi tiết.
\item Xây dựng phần trắc nghiệm ôn tập với phản hồi đúng, sai, giải thích đáp án và kết quả làm bài để hỗ trợ người học tự kiểm tra kiến thức.
\item Tổng hợp vai trò lãnh đạo của Đảng, nguyên nhân thắng lợi, ý nghĩa lịch sử và bài học kinh nghiệm để giúp người học liên kết kiến thức giữa các giai đoạn.
\item Tổ chức giao diện và điều hướng thống nhất giữa các khu vực, bảo đảm mối liên hệ giữa hoạt động đọc học liệu, xem tư liệu và ôn tập.
\end{enumerate}

\section{Đối tượng sử dụng và phạm vi sản phẩm}

\textbf{Đối tượng sử dụng chủ yếu} là sinh viên học môn Lịch sử Đảng Cộng sản Việt Nam, đặc biệt những người cần hệ thống hóa kiến thức về giai đoạn 1954 - 1975. Website cũng được định hướng làm phương tiện minh họa khi thảo luận hoặc trình bày nội dung học tập theo nhóm.

\textbf{Về nội dung}, học liệu tập trung vào sự lãnh đạo của Đảng trong kháng chiến chống Mỹ, cứu nước, gắn nhiệm vụ xây dựng miền Bắc với đấu tranh giải phóng miền Nam. Các sự kiện được lựa chọn nhằm làm rõ mối quan hệ giữa bối cảnh, đường lối và thực tiễn, không đặt mục tiêu tái hiện toàn bộ diễn biến chiến tranh. Việc phân chia thành năm giai đoạn phục vụ tổ chức kiến thức và điều hướng trên website.

\textbf{Về nguồn tài liệu}, sản phẩm sử dụng \emph{Giáo trình Lịch sử Đảng Cộng sản Việt Nam} do Bộ Giáo dục và Đào tạo tổ chức biên soạn, Nhà xuất bản Chính trị quốc gia Sự thật xuất bản năm 2021 làm cơ sở học thuật chính. Trọng tâm là phần lãnh đạo xây dựng chủ nghĩa xã hội ở miền Bắc và kháng chiến chống Mỹ, cứu nước trong Chương 2. Tư liệu từ Bảo tàng Chứng tích Chiến tranh được sử dụng để liên hệ, minh họa cho nội dung lịch sử. [1, tr. 180 - 235].

\textbf{Về chức năng}, website tập trung vào ba hoạt động: tìm hiểu kiến thức qua trình chiếu và dòng thời gian; xem tư liệu chứng tích theo giai đoạn; tự ôn tập bằng trắc nghiệm có giải thích. Phần tổng kết và bài học hỗ trợ khái quát nội dung, còn phần giới thiệu nhóm cung cấp thông tin về tập thể thực hiện sản phẩm.

\textbf{Về giá trị sử dụng}, sản phẩm được định hướng hỗ trợ tự học và trình bày kiến thức trong phạm vi học phần. Kết quả trắc nghiệm phản ánh mức độ trả lời đúng đối với bộ câu hỏi của website, không thay thế đánh giá chính thức của giảng viên.

\section{Phương pháp thực hiện}

\textbf{Nghiên cứu và chọn lọc nội dung} được tiến hành trên cơ sở giáo trình, kết hợp phương pháp lịch sử và phương pháp logic. Các sự kiện được đặt theo trình tự thời gian và trong hoàn cảnh cụ thể; đồng thời làm rõ quan hệ giữa mục tiêu, chủ trương, tổ chức thực hiện và kết quả. Đây là cơ sở để lựa chọn kiến thức phù hợp với chủ đề của website.

\textbf{Biên soạn học liệu số} được thực hiện bằng cách phân tích, tổng hợp và chuyển nội dung nghiên cứu thành những đơn vị kiến thức ngắn gọn. Các giai đoạn sử dụng chung một cấu trúc trình bày để thuận tiện theo dõi và so sánh. Nội dung tổng kết kết nối các giai đoạn, trong khi câu hỏi ôn tập tập trung vào những luận điểm và sự kiện trọng tâm.

\textbf{Khai thác tư liệu trực quan} kết hợp hiện vật, ảnh tư liệu, nội dung giới thiệu của bảo tàng và ghi nhận khảo sát của các thành viên. Tư liệu được xem xét theo nguồn gốc, bối cảnh và nội dung phản ánh trước khi liên hệ với kiến thức lịch sử. Cách thực hiện này giúp hình ảnh giữ vai trò bổ trợ cho nhận thức, thay vì chỉ làm yếu tố trang trí.

\textbf{Tổ chức trải nghiệm học tập trên website} dựa trên mối liên hệ giữa đọc, quan sát và tự kiểm tra. Dòng thời gian hỗ trợ định vị giai đoạn; khu vực trình chiếu trình bày kiến thức; kho tư liệu bổ sung hình ảnh và thông tin; phần trắc nghiệm đưa ra phản hồi cùng giải thích. Các khu vực được kết nối bằng hệ thống điều hướng để người học chủ động lựa chọn nội dung phù hợp với nhu cầu.

\textbf{Rà soát sản phẩm} tập trung vào sự nhất quán giữa nội dung lịch sử, chú thích tư liệu, câu hỏi và đáp án, cùng cách hiển thị và điều hướng giữa các khu vực. Tiêu chí xem xét là mức độ rõ ràng của học liệu, khả năng thực hiện các thao tác cơ bản và sự phù hợp với mục tiêu hỗ trợ học tập của môn học.

\clearpage
